Let $L$ = luminosity of the giant stars in the giant branch,
$L_{MS}$ = luminosity of the main sequence star,
$F$ = flux of the giant stars in the giant branch,
$F_{MS}$ = flux of the main sequence star,
$T$ = temperature of the giant star,
$T_{MS}$ = temperature of the main sequence star,
$R$ = radius of the giant star,
$R_{MS}$ = radius of the main sequence star.

Use the radius--luminosity--temperature relation
\[
  \frac{L}{L_{MS}} = \frac{R^2 T^4}{R_{MS}^2 T_{MS}^4}
    = \left(\frac{R}{R_{MS}}\right)^{\!2}\left(\frac{T}{T_{MS}}\right)^{\!4}
    = \left(100\right)^2 \left(\frac{1}{3}\right)^{\!4} = 123.5
\]
\hfill(50\%)
\begin{align*}
  F &= \frac{L}{4\pi d^2} \\
  \frac{F}{F_{MS}} &= \frac{L}{L_{MS}}\left(\frac{d_{MS}}{d}\right)^{\!2}
\end{align*}
that is, the star becomes 123.5 times brighter when it becomes a giant star. To
be just barely visible using the same telescope, we push the star farther so
that it becomes 123.5 times dimmer (flux is inversely proportional to a square
of the distance).
\begin{align*}
  \frac{1}{123.5} &= \left[\frac{1/d}{1/d_{MS}}\right]^{2}
  \;\Longrightarrow\; \frac{d}{d_{MS}} = 11.11
  \;\Longrightarrow\; d = 11.11\,d_{MS} \\
  &= 11.11\,(20)\ \text{pc} = 222.20\ \text{pc}
\end{align*}
\hfill(50\%)
